% GENERATED FILE. Edit canonical lessons, then run npm run generate:books.
% canonical-source-hash: fnv1a64:783e5b7c43b01abd
% canonical-lessons: ES-C695-picar, ES-C695-fauna, ES-C695-flora, ES-C695-pescar, ES-C695-ciervo, ES-R695-repaso-688-695

\chapter{Fauna y flora}
\label{ch:es-fauna-y-flora}
\hlchaptermodality{695}

\begin{chapteropening}
\textbf{By the end of this chapter:} \emph{I can talk about wildlife such as deer, plant life, fishing, and insects that bite.}

Complete the last lesson of chapter 695: I can talk about wildlife such as deer, plant life, fishing, and insects that bite.
\end{chapteropening}

\section[picar]{picar — to bite, to sting (insects); to be spicy; to nibble}

\label{lesson:ES-C695-picar}



\begin{quote}
Before the new one: say the Spanish for the butterfly, then the Spanish for the small lizard.
\end{quote}

\subsection*{You'll want to know: picar}
\textbf{picar} — \textquotedblleft{}to bite, to sting (insects); to be spicy; to nibble\textquotedblright{}.

\begin{quote}
Me picó un mosquito en el brazo.
\end{quote}

\begin{quote}
\emph{A mosquito bit me on the arm.}
\end{quote}

\emph{Esta salsa pica} — this sauce is hot. \emph{Picar algo} — to have a snack.

\begin{cousinweb}
\emph{picar} comes from Vulgar Latin \emph{piccāre} \textquotedblleft{}to prick, to strike\textquotedblright{} (EtymDB); where that verb came from is not settled.
\end{cousinweb}

\begin{grammarlens}[title={What you've built}]
One more word for telling, explaining and describing at B1.
\end{grammarlens}

\subsection*{Your turn}
\begin{itemize}
\raggedright
  \item \emph{Say it:} \emph{picar}
  \item \emph{Say it:} \emph{picar}, once more
  \item \emph{Say it:} a sentence of your own with \emph{picar}
  \item \emph{Recall:} say the Spanish for the butterfly, then the Spanish for the small lizard, then say \emph{picar} again
\end{itemize}

\begin{tcolorbox}[breakable,colback=teal!4,colframe=teal!35!black,title={Before you move on}]
What does picar mean? (\textquotedblleft{}To bite, to sting (insects); to be spicy; to nibble\textquotedblright{}.) Say it once more.
\end{tcolorbox}
\section[la fauna]{la fauna — the fauna, the wildlife}

\label{lesson:ES-C695-fauna}



\begin{quote}
Before the new one: say the Spanish for the small lizard, then the Spanish for to bite.
\end{quote}

\subsection*{You'll want to know: la fauna}
\textbf{la fauna} — \textquotedblleft{}the fauna, the wildlife\textquotedblright{}.

\begin{quote}
La fauna de la isla es única.
\end{quote}

\begin{quote}
\emph{The wildlife on the island is unique.}
\end{quote}

The animals of a place: \emph{la fauna marina} — sea life.

\begin{cousinweb}
\emph{fauna} is a learned borrowing from Latin \emph{Fauna}, a Roman goddess, the sister of the woodland god \emph{Faunus} (Lewis \& Short).
\end{cousinweb}

\begin{grammarlens}[title={What you've built}]
One more word for telling, explaining and describing at B1.
\end{grammarlens}

\subsection*{Your turn}
\begin{itemize}
\raggedright
  \item \emph{Say it:} \emph{la fauna}
  \item \emph{Say it:} \emph{la fauna}, once more
  \item \emph{Say it:} a sentence of your own with \emph{la fauna}
  \item \emph{Recall:} say the Spanish for the small lizard, then the Spanish for to bite, then say \emph{la fauna} again
\end{itemize}

\begin{tcolorbox}[breakable,colback=teal!4,colframe=teal!35!black,title={Before you move on}]
What does la fauna mean? (\textquotedblleft{}The fauna, the wildlife\textquotedblright{}.) Say it once more.
\end{tcolorbox}
\section[la flora]{la flora — the flora, the plant life}

\label{lesson:ES-C695-flora}



\begin{quote}
Before the new one: say the Spanish for to bite, then the Spanish for the fauna.
\end{quote}

\subsection*{You'll want to know: la flora}
\textbf{la flora} — \textquotedblleft{}the flora, the plant life\textquotedblright{}.

\begin{quote}
Aquí la flora es muy variada.
\end{quote}

\begin{quote}
\emph{The plant life here is very varied.}
\end{quote}

Often with \emph{la fauna}.

\begin{cousinweb}
\emph{flora} is a learned borrowing from Latin \emph{Flōra}, the Roman goddess of flowers (Lewis \& Short), from \emph{flōs} \textquotedblleft{}flower\textquotedblright{}.
\end{cousinweb}

\begin{grammarlens}[title={What you've built}]
One more word for telling, explaining and describing at B1.
\end{grammarlens}

\subsection*{Your turn}
\begin{itemize}
\raggedright
  \item \emph{Say it:} \emph{la flora}
  \item \emph{Say it:} \emph{la flora}, once more
  \item \emph{Say it:} a sentence of your own with \emph{la flora}
  \item \emph{Recall:} say the Spanish for to bite, then the Spanish for the fauna, then say \emph{la flora} again
\end{itemize}

\begin{tcolorbox}[breakable,colback=teal!4,colframe=teal!35!black,title={Before you move on}]
What does la flora mean? (\textquotedblleft{}The flora, the plant life\textquotedblright{}.) Say it once more.
\end{tcolorbox}
\section[pescar]{pescar — to fish; to catch}

\label{lesson:ES-C695-pescar}



\begin{quote}
Before the new one: say the Spanish for the fauna, then the Spanish for the flora.
\end{quote}

\subsection*{You'll want to know: pescar}
\textbf{pescar} — \textquotedblleft{}to fish; to catch\textquotedblright{}.

\begin{quote}
Mi padre va a pescar al río.
\end{quote}

\begin{quote}
\emph{My father goes fishing in the river.}
\end{quote}

\emph{El pescado} is fish as food; \emph{el pez} is a live fish.

\begin{cousinweb}
\emph{pescar} is inherited from Vulgar Latin \emph{piscāre}, Latin \emph{piscārī} \textquotedblleft{}to fish\textquotedblright{} (EtymDB), from \emph{piscis} \textquotedblleft{}fish\textquotedblright{}.
\end{cousinweb}

\begin{grammarlens}[title={What you've built}]
One more word for telling, explaining and describing at B1.
\end{grammarlens}

\subsection*{Your turn}
\begin{itemize}
\raggedright
  \item \emph{Say it:} \emph{pescar}
  \item \emph{Say it:} \emph{pescar}, once more
  \item \emph{Say it:} a sentence of your own with \emph{pescar}
  \item \emph{Recall:} say the Spanish for the fauna, then the Spanish for the flora, then say \emph{pescar} again
\end{itemize}

\begin{tcolorbox}[breakable,colback=teal!4,colframe=teal!35!black,title={Before you move on}]
What does pescar mean? (\textquotedblleft{}To fish; to catch\textquotedblright{}.) Say it once more.
\end{tcolorbox}
\section[el ciervo]{el ciervo — the deer, the stag}

\label{lesson:ES-C695-ciervo}



\begin{quote}
Before the new one: say the Spanish for the flora, then the Spanish for to fish.
\end{quote}

\subsection*{You'll want to know: el ciervo}
\textbf{el ciervo} — \textquotedblleft{}the deer, the stag\textquotedblright{}.

\begin{quote}
Por la mañana vimos un ciervo en el bosque.
\end{quote}

\begin{quote}
\emph{In the morning we saw a deer in the woods.}
\end{quote}

The female is \emph{la cierva}.

\begin{cousinweb}
\emph{ciervo} is inherited from Latin \emph{cervus} \textquotedblleft{}a stag, a deer\textquotedblright{} (Lewis \& Short).
\end{cousinweb}

\begin{grammarlens}[title={What you've built}]
One more word for telling, explaining and describing at B1.
\end{grammarlens}

\subsection*{Your turn}
\begin{itemize}
\raggedright
  \item \emph{Say it:} \emph{el ciervo}
  \item \emph{Say it:} \emph{el ciervo}, once more
  \item \emph{Say it:} a sentence of your own with \emph{el ciervo}
  \item \emph{Recall:} say the Spanish for the flora, then the Spanish for to fish, then say \emph{el ciervo} again
\end{itemize}

\begin{tcolorbox}[breakable,colback=teal!4,colframe=teal!35!black,title={Before you move on}]
What does el ciervo mean? (\textquotedblleft{}The deer, the stag\textquotedblright{}.) Say it once more.
\end{tcolorbox}
\section[(dialogue)]{Repaso de los capítulos 688–695}

\label{lesson:ES-R695-repaso-688-695}



\begin{quote}
Say the Spanish for the ocean and the wave.
\end{quote}

\subsection*{Your turn}
Say these aloud:

\begin{itemize}
\raggedright
  \item the ocean, the wave, the tide, the bay, to float
  \item the waterfall, the stream, the lagoon, the marsh, the cliff
  \item to have a swim, to get a tan, the sunshade, the swimsuit, to lie down
  \item the farmer, the country person, the country estate, the vineyard, to grow
  \item to sow, to water, to harvest, the well, to flower
  \item the olive tree, the pine tree, the bush, the vegetation, the oak tree
  \item the seagull, the large bird, the insect, the butterfly, the small lizard
  \item to bite, the fauna, the flora, to fish, the deer
\end{itemize}

\begin{tcolorbox}[breakable,colback=teal!4,colframe=teal!35!black,title={Before you move on}]
The ocean? (\textbf{el océano}.) The wave? (\textbf{la ola}.) The tide? (\textbf{la marea}.) The bay? (\textbf{la bahía}.) To float? (\textbf{flotar}.) The waterfall? (\textbf{la cascada}.) The stream? (\textbf{el arroyo}.) The lagoon? (\textbf{la laguna}.) The marsh? (\textbf{el pantano}.) The cliff? (\textbf{el acantilado}.) To have a swim? (\textbf{bañarse}.) To get a tan? (\textbf{broncearse}.) The sunshade? (\textbf{la sombrilla}.) The swimsuit? (\textbf{el bañador}.) To lie down? (\textbf{tumbarse}.) The farmer? (\textbf{el agricultor}.) The country person? (\textbf{el campesino}.) The country estate? (\textbf{la finca}.) The vineyard? (\textbf{el viñedo}.) To grow? (\textbf{cultivar}.) To sow? (\textbf{sembrar}.) To water? (\textbf{regar}.) To harvest? (\textbf{cosechar}.) The well? (\textbf{el pozo}.) To flower? (\textbf{florecer}.) The olive tree? (\textbf{el olivo}.) The pine tree? (\textbf{el pino}.) The bush? (\textbf{el arbusto}.) The vegetation? (\textbf{la vegetación}.) The oak tree? (\textbf{el roble}.) The seagull? (\textbf{la gaviota}.) The large bird? (\textbf{el ave}.) The insect? (\textbf{el insecto}.) The butterfly? (\textbf{la mariposa}.) The small lizard? (\textbf{la lagartija}.) To bite? (\textbf{picar}.) The fauna? (\textbf{la fauna}.) The flora? (\textbf{la flora}.) To fish? (\textbf{pescar}.) The deer? (\textbf{el ciervo}.)
\end{tcolorbox}
